\chapter*{Предисловие}
\addcontentsline{toc}{chapter}{Предисловие}
\markboth{Предисловие}{Предисловие}

Эта книга отвечает на вопрос, который сначала казался простым: \emph{как дрон LiteBee Wing понимает, над какой клеткой ковра он находится}. Оказалось, что за ответом стоит целая цепочка. Камера на дроне, отдельный маленький компьютер под Linux, программа распознавания меток, кабель с собственным форматом кадра, автопилот со своей формулой и приложение, которое всё это показывает. Мы прошли эту цепочку от начала до конца и записали всё, что удалось подтвердить.

\section*{Как читать книгу}

Главы идут от общей картины к деталям. Если нужна только суть, читайте главы~\ref{ch:kartina}, \ref{ch:uart} и~\ref{ch:kovyor}. Если интересует, как устроена система и сеть, читайте главу~\ref{ch:sys}. Если хотите писать свою прошивку или свою камеру, главы~\ref{ch:svoya} и~\ref{ch:pyatna} скажут, что уже известно и где ещё тёмные пятна.

\begin{zametka}[Пометки достоверности]
Вся книга построена на трёх пометках, и мы стараемся не путать их.
\begin{itemize}
  \item \pr{} --- факт виден прямо в файле, в дизассемблере, в памяти работающей программы или в ответе устройства.
  \item \vyv{} --- наше заключение из прочитанного. Оно правдоподобно, но не проверено напрямую.
  \item \nei{} --- мы не знаем. Это честнее, чем гадать.
\end{itemize}
\end{zametka}

\section*{Что у нас было}

\begin{itemize}
  \item Образы прошивки дрона Wing от производителя и пакеты обновления (Linux на Allwinner A33).
  \item Прошивка ядра полёта C10 (RISC-V) с отладочной информацией: имена функций и структур.
  \item Приложение LiteBee Go для компьютера и два мобильных приложения.
  \item \textbf{Сам дрон} и \textbf{сам модуль с камерой}, подключённые к компьютеру по USB.
\end{itemize}

\begin{vnimanie}[Что нужно помнить]
Дрон, на котором мы работали, использует прошивку марта 2023 года, а образ, который прислал производитель, новее (версия 1.2.8, 2025 год). Их автопилоты отличаются. Везде, где это важно, мы указываем, о какой версии идёт речь. Формула и формат кадра между ними совпали, а вот заголовок кадра отличается (глава~\ref{ch:uart}).
\end{vnimanie}

\begin{opasno}
Все эксперименты с дроном выполнялись \textbf{без пропеллеров} и только на чтение: ни во флеш дрона, ни во флеш модуля ничего не записывалось. Если вы повторяете что-то из книги, делайте так же.
\end{opasno}
